\section{Conclusions}
\label{sec:conclusions}

This paper introduced CascadeBench, a black-box benchmark that measures how
acoustic channel degradations propagate through commercial speech
recognition and translation services, and reported a complete measurement
over a controlled bilingual corpus of 480 sentences, fourteen degradation
conditions from seven families, and an eight-level noise severity sweep
from 20 to $-10$~dB, covering five recognition
systems from three product classes, three cascade translation paths,
and two end-to-end speech translation paths, all accessed through their
public interfaces on free service quotas, and validated on 150
human-recorded utterances with measured impulse responses and real noise,
with three hold-out measurements --- an open-source locally run recogniser,
a family-disjoint third translation reference and backend, and a
measured-RIR reverb calibration --- probing how far the findings depend on
the single vendor that produced the systems, the references, and the
synthesis.

The measurements support five conclusions. First, the product classes of
commercial recognisers differ in where they fail under noise, not merely how
fast: LLM-based recognisers hold a sub-percent floor and cross ten percent
error 2.2 to 2.5~dB later in Mandarin and 4.4 to 4.8~dB later in English
than the specialised streaming system, then collapse more sharply, and below
0~dB they fail by insertion-heavy hallucination with character error rates
above 100 percent; a replication on 150 human-recorded utterances with
measured impulse responses and real multi-talker babble reproduces this
knee-and-collapse structure and the hallucination regime, with two
calibrations: real rooms at matched reverberation time prove far milder
than the synthetic low-DRR stress response, and the LLM noise-tolerance
advantage holds only in the near-knee, low-error regime, disappearing and
for one LLM system reversing once error passes the ten-percent level into
the hallucination regime. Second, robustness is
condition- and language-specific
rather than uniform: narrowbanding, coding, packet loss, clipping, and tempo
changes at deployed severities are individually negligible for all systems,
0~dB babble nearly destroys English recognition for every system while
leaving the best Mandarin recognition low, and no system, of any class,
transcribes reliably below 0~dB. Third, the quality of a speech translation
cascade under degradation is governed by its recognition frontend, and the
relationship is an empirical linear pass-through of 0.44 to 0.54
BLEU points per percentage point of recognition error (condition-level 95
percent intervals of 0.23--0.45, 0.52--0.66, and a heavily right-tailed
0.38--3.9 for the three cascades), with intercepts
statistically indistinguishable from zero, stable under leave-one-condition-out
refits for two of the three cascades, and honest condition-level
uncertainty that is substantially wider than the sentence-level intervals,
across up to fifteen conditions and three frontends; on real speech the
same relation holds with slopes of 0.41 and 0.73 for the two translation
directions. Fourth,
end-to-end speech translation does not inherit
robustness automatically: the most resilient translation path measured on
synthetic speech is a
dedicated streaming end-to-end service, while the omni-modal end-to-end
path is both the slowest and, under severe noise, the most fragile on the
point estimate, and on real
speech the streaming end-to-end path loses a larger fraction of its clean
quality under 0~dB real babble than a well-frontended cascade (a paired
contrast of $+7.9$ points whose 95 percent interval still spans zero),
while the
absolute BLEU ordering of the two paths remains reference-dependent. Fifth,
black-box evaluation carries specific methodological obligations that this
benchmark operationalises, including multi-reference scoring against at
least two model families to control machine-reference affinity (here the
relative-loss ordering proved stable across same-vendor, oracle, and
family-disjoint GLM references, while magnitudes moved by up to eleven
points per system), exact-transcript corpora to enable sub-percent
measurement, the reporting of response completeness alongside error rates,
and the anchoring of service-level findings in measured signal properties,
which here exposed a sharp dissociation between intelligibility metrics and
recognition error across degradation families. The three hold-out
measurements bound the single-vendor caveat rather than removing it: the
knee-and-collapse structure, the language-asymmetric hallucination regime,
and the cascade relative-loss band all replicate with no vendor component
in common, but the absolute numbers --- the knee locations, the slopes, the
score levels --- remain measurements of one vendor's services at one access
date, and the open-source recogniser used for the replication is a
small, CPU-quantised model whose absolute accuracy is, as expected, well
below the commercial systems it mirrors structurally.

The benchmark recipe, including the corpus, the seeded degradation pipeline,
the degraded audio tree, all model responses, and the analysis code, is
released so that the measurement can be re-executed and verified as
providers update their services. The released artefacts are versioned alongside the service
identifiers and access dates of every system, so that future re-runs can
attribute score movements either to provider-side model updates or to
changes in the benchmark itself, a distinction that periodic evaluation of
moving targets requires. Together with the zero-cost execution protocol
described in Section~\ref{sec:methods}, the release is intended to make
robustness re-measurement a routine maintenance activity for speech
deployments rather than a one-off research exercise. As
speech intelligence continues to migrate into provider-operated black-box
systems, controlled measurement of this kind, repeated over time and across
providers, offers the community a practical substitute for the training-set
access that classical robustness research assumed, and offers integrators
the quantitative basis that architectural intuition alone cannot supply.
